\documentclass[11pt]{article}
\usepackage[margin=1in]{geometry}
\usepackage{amsmath}
\usepackage{booktabs}
\usepackage{hyperref}
\emergencystretch=2em

\title{Minimal Symbolic Pipelines for Constraint-Checked State Records}
\author{Lijie Wang\\Independent Researcher}
\date{May 2026}

\begin{document}
\maketitle

\begin{abstract}
This technical note describes a clean-room minimal research artifact for
evaluating a constraint-checked symbolic pipeline over normalized records. The
artifact follows a fixed implementation sequence:
\[
\text{SyntheticCase} \rightarrow \text{TraceStore} \rightarrow
\text{Baselines} \rightarrow \text{ConstraintCheck} \rightarrow
\text{ScoreCard}.
\]
The goal is to make the separation between structured records, intermediate
trace state, method proposals, explicit constraint checks, and aggregate scores
directly inspectable. The artifact does not include natural language parsing,
persona generation, user profiling, product-specific memory logic, or
production code. It is intended as an external-review draft of a small
artifact paper / technical note, not as a broad empirical claim.
\end{abstract}

\section{Introduction}

Constraint-checked state records are useful when a pipeline must distinguish
between what is represented, what a method proposes, what a constraint checker
accepts, and what an aggregate score reports. In small multi-stage symbolic
systems, these roles can easily be collapsed into a single output value. Once
that happens, it becomes difficult to inspect whether a result came from the
record representation, the candidate-selection method, the constraint layer, or
the score aggregation rule.

This paper studies that separation in a deliberately small setting. The
artifact uses normalized symbolic records rather than free-form text. Each case
contains structured traces and finite lists of allowed, denied, and required
state labels. A deterministic method produces a candidate record, a constraint
checker evaluates that candidate against the case lists, and a score card
aggregates the checked records.

The scope is intentionally narrow. This work does not make claims about
general language models, persona models, user modeling, or deployed systems.
It studies a compact reproducible artifact whose purpose is to make the
interfaces between records, traces, proposals, checks, and scores explicit.

The contribution is an inspectable artifact consisting of:
\begin{itemize}
    \item a fixed bank of 20 synthetic normalized cases;
    \item three deterministic baseline methods;
    \item an explicit finite constraint checker; and
    \item a reproducible CSV score summary generated from checked records.
\end{itemize}

\section{Related Work and Positioning}

This artifact is closest in spirit to small reproducible symbolic evaluation
pipelines rather than learned language models. The constraint-checking layer is
related to the long tradition of constraint satisfaction and consistency
checking, where explicit relations are manipulated and tested rather than
learned from examples \cite{mackworth1977}. The state-construction part is also
adjacent to program-analysis work in which concrete records are mapped into
intermediate abstract states for inspection or approximation
\cite{cousot1977}. The present artifact does not attempt to reproduce the
generality of those frameworks. Instead, it borrows the narrower idea that
state, candidate, and checking interfaces should be explicit.

The repository is also positioned as a reproducible computational artifact.
Reproducible research has emphasized the importance of shipping enough data,
code, and scripts for a reader to regenerate reported computational results
\cite{claerbout1992}. The artifact follows this convention at a small scale:
the case bank, deterministic baselines, pilot command, tests, and CSV summary
are included so that the reported table can be regenerated directly.

This positioning is intentionally modest. The work is not a new constraint
solver, a static analyzer, or a learned model. It is a compact technical note
that isolates one checked-record interface for external review.

\section{Problem Formulation}

The paper-level formulation is the following checked-record chain:
\[
r_t \rightarrow s_t \rightarrow p_t \rightarrow z_t \rightarrow q.
\]

Here $r_t$ is a normalized input record, $s_t$ is the trace state constructed
from observed records, $p_t$ is a method proposal or candidate, $z_t$ is the
checked record after applying constraints, and $q$ is an aggregate score card.
In the implemented artifact, the transformation can be written as:
\[
s_t = F(s_{t-1}, r_t),
\]
\[
p_t = G(s_t, m),
\]
\[
z_t = H(p_t, C),
\]
\[
q = A(\{z_t\}).
\]
The symbol $m$ denotes the candidate-selection method, and $C$ denotes the
finite constraint set attached to a case.

This formulation does not require natural language parsing or a learned model.
It requires only structured input records, deterministic candidate methods, and
finite reproducible constraints. The design is therefore small enough to audit
directly: each checked output record can be traced back to an input case, an
intermediate state, a method choice, and an explicit constraint decision.

\section{Artifact Design}

The implementation-level artifact follows this sequence:
\[
\text{SyntheticCase} \rightarrow \text{TraceStore} \rightarrow
\text{Baselines} \rightarrow \text{ConstraintCheck} \rightarrow
\text{ScoreCard}.
\]
This sequence is the concrete module boundary for the repository. It is the
implementation counterpart to the symbolic chain in the previous section.

\subsection{SyntheticCase}

The case bank contains fixed synthetic normalized records. Each case has a
case identifier, a list of traces, an allow list, a deny list, and a required
match list. Each trace has a trace identifier, a signal label, two symbolic
agent identifiers, a tick value, and structured attributes. The attributes
include a state label, an origin trace, a weight band, and a decay rate. The
case bank is synthetic and is included only to make the protocol reproducible.

\subsection{TraceStore}

The trace store provides the shared intermediate representation. It orders
case traces by tick and trace identifier, then exposes them as records for
state construction. This stage does not select a method-specific candidate. Its
role is to provide a common ordered trace input to all baselines.

\subsection{Baselines}

The artifact includes three deterministic baseline methods:
\texttt{summary\_loop}, \texttt{template\_grid}, and \texttt{symbolic\_rule}.
All three methods receive the same constructed trace state and differ only in
candidate selection.

In method-type terms, \texttt{summary\_loop} is a max-score baseline: it
selects the highest-scoring accumulated state label. \texttt{template\_grid} is
a position-index baseline: it selects a state label by using the record count
and pass index as a deterministic position rule. \texttt{symbolic\_rule} is an
allow-filtered symbolic baseline: it first filters candidate labels through the
case allow list when possible, then selects among the remaining accumulated
state labels. These baselines are simple by design; their purpose is to make
method-level differences visible under the same checked-record protocol.

\subsection{ConstraintCheck}

The constraint checker evaluates a proposal against explicit finite
constraints. For each checked record, it reports whether the selected state
label appears in the allow list, whether it appears in the deny list, and
whether it appears in the required match list. A checked record passes only
when the selected state is allowed, not denied, and required by the case.

\subsection{ScoreCard}

The score card aggregates checked records by method. It reports the number of
runs, average score, pass rate, and average record count. These metrics are not
intended to rank methods in general. They summarize behavior under the fixed
case bank and constraint protocol.

\section{Pilot Protocol}

The pilot uses 20 cases, three methods, and ten repeated passes:
\[
20 \times 3 \times 10 = 600
\]
checked records. The exact command used to generate the pilot records is:
\begin{verbatim}
python -m scripts.run_pilot --cases data/cases_small.jsonl --runs 10 \
  --methods summary_loop,template_grid,symbolic_rule \
  --out outputs/pilot_runs.jsonl
\end{verbatim}
The aggregate CSV is generated by:
\begin{verbatim}
python -m scripts.export_summary --input outputs/pilot_runs.jsonl \
  --out outputs/summary_agg.csv
\end{verbatim}

The repeated passes should not be interpreted as estimating randomness. The
methods are deterministic. The pass index is part of the protocol because one
baseline uses it as a deterministic selection position, and because repeated
passes produce a complete checked-record table for the fixed protocol. The
pilot therefore verifies the record-generation protocol and exposes method
behavior across the same case bank, rather than estimating stochastic
performance. No statistical tests are applied because the protocol is
deterministic and fully enumerated for the fixed case bank.

\section{Results}

Table~\ref{tab:results} reports the aggregate pilot metrics generated from the
current artifact. The pilot contains 600 checked records, grouped into 200
records per method.

\begin{table}[htbp]
\centering
\caption{Aggregate pilot metrics generated from the checked-record protocol.}
\label{tab:results}
\begin{tabular}{lccc}
\toprule
Method & Average score & Pass rate & Average record count \\
\midrule
summary\_loop & 1.392134 & 0.100000 & 6.850000 \\
template\_grid & 0.793868 & 0.180000 & 6.850000 \\
symbolic\_rule & 0.917717 & 0.350000 & 6.850000 \\
\bottomrule
\end{tabular}
\end{table}

The pilot exposes a score/pass-rate tradeoff. The \texttt{summary\_loop}
baseline has the highest average score, $1.392134$, but the lowest pass rate,
$0.100000$. The \texttt{symbolic\_rule} baseline has a lower average score,
$0.917717$, but the highest pass rate, $0.350000$. The
\texttt{template\_grid} baseline has the lowest average score, $0.793868$, and
an intermediate pass rate, $0.180000$.

These results show that the score and constraint satisfaction are different
dimensions in this artifact. Selecting the highest accumulated score does not
necessarily select a candidate that satisfies the finite case constraints. A
constraint-aware selection rule can increase the pass rate under this protocol
while producing a lower average score. This observation is specific to the
fixed synthetic case bank and deterministic baselines; it should be read as a
pilot comparison, not as a general method ranking.

Table~\ref{tab:failures} reports a constraint failure breakdown for the same
checked records. The failure columns are not mutually exclusive: a single
checked record may fail the allow condition and the required-match condition at
the same time.

\begin{table}[htbp]
\centering
\caption{Constraint failure breakdown over 200 checked records per method.
``Deny hit'' means that the selected label appears in the deny list; failure
columns are not mutually exclusive.}
\label{tab:failures}
\begin{tabular}{lrrrr}
\toprule
Method & Allow failed & Deny hit & Required failed & Passed \\
\midrule
summary\_loop & 150 & 30 & 180 & 20 \\
template\_grid & 116 & 27 & 164 & 36 \\
symbolic\_rule & 0 & 0 & 130 & 70 \\
\bottomrule
\end{tabular}
\end{table}

The breakdown explains the aggregate tradeoff more directly. The
\texttt{summary\_loop} baseline often selects high-scoring labels that do not
match the case constraints. The \texttt{symbolic\_rule} baseline avoids allow
and deny failures in this pilot, but still fails the required-match condition
in 130 of 200 checked records. This supports the central artifact claim: the
candidate score, allow/deny checks, required-match checks, and final pass rate
are separable quantities that should be reported separately.

\section{Limitations}

This artifact has several important limitations. The case bank is synthetic and
contains only 20 cases. The baselines are deterministic and intentionally
simple. The constraints are finite and manually specified within each case. The
artifact does not include learned methods, natural language parsing, or any
claim about broad empirical behavior. It also does not include production
logic or a deployment claim.

These limitations are part of the design boundary. The purpose of the artifact
is to provide a small reproducible object for external review, not to establish
large-scale performance claims.

\section{Conclusion}

This technical note presents a minimal reproducible basis for studying
constraint-checked state records. The contribution is not performance. The
contribution is a clean separation between normalized records, shared trace
state, method proposals, explicit constraint checks, and aggregate scores.

Future work should stay within this interface while increasing the strength of
the study. Natural next steps include stronger baselines, larger symbolic case
banks, and formal analysis of constraint-preserving state transitions. These
extensions can be evaluated against the same basic record-to-check-to-score
boundary introduced here.

\begin{thebibliography}{9}

\bibitem{mackworth1977}
Alan K. Mackworth.
\newblock Consistency in networks of relations.
\newblock \emph{Artificial Intelligence}, 8(1):99--118, 1977.

\bibitem{cousot1977}
Patrick Cousot and Radhia Cousot.
\newblock Abstract interpretation: a unified lattice model for static analysis
of programs by construction or approximation of fixpoints.
\newblock In \emph{Conference Record of the Sixth Annual ACM SIGPLAN-SIGACT
Symposium on Principles of Programming Languages}, pages 238--252, 1977.

\bibitem{claerbout1992}
Jon F. Claerbout and Martin Karrenbach.
\newblock Electronic documents give reproducible research a new meaning.
\newblock In \emph{SEG Technical Program Expanded Abstracts}, pages 601--604,
1992.

\end{thebibliography}

\end{document}
